\chapter{Lạm phát, cấu trúc, và sự mất đi của các điều kiện ban đầu}

\section{Omega và bằng chứng về độ phẳng}

Một câu hỏi của sinh viên về \(\Omega\) lớn khiến cuộc thảo luận quay trở lại phương trình Friedmann-Robertson-Walker. Đối với kỷ nguyên hiện nay, khá lâu sau thời kỳ bức xạ chi phối, định luật giãn nở được viết với một hạng năng lượng chân không và một hạng vật chất, còn bức xạ thì được bỏ qua vì hiện nay nó rất nhỏ:
\begin{equation}
H^2=\left(\frac{\dot a}{a}\right)^2=\frac{8\pi G}{3}\left(\rho_0+\rho_m\right)-\frac{k}{a^2}.
\end{equation}
Hạng chân không không bị pha loãng. Hạng vật chất thì có:
\begin{equation}
\rho_m=\frac{M}{a^3}.
\end{equation}
\(\Omega\) là phần đóng góp của các thành phần khác nhau vào phương trình giãn nở.

Chia cho \(H^2\). Khi đó các hạng đã chuẩn hóa cộng lại bằng một. Trong một ký hiệu đã được chuẩn hóa,
\begin{equation}
\begin{aligned}
\Omega_0&=\frac{8\pi G\rho_0}{3H^2},\\
\Omega_m&=\frac{8\pi G\rho_m}{3H^2},\\
\Omega_k&=-\frac{k}{a^2H^2},
\end{aligned}
\end{equation}
và do đó
\begin{equation}
\Omega_0+\Omega_m+\Omega_k=1.
\end{equation}
\editorialnote{Âm thanh và phần bảng ở đoạn này khá khó theo dõi, và người giảng chuyển qua lại giữa một vài tên gọi cho phần đóng góp chân không. Ký hiệu ở đây chuẩn hóa lại lập luận sau khi chia cho \(H^2\).}
Nếu không gian phẳng, \(k=0\), nên
\begin{equation}
\Omega_0+\Omega_m=1.
\end{equation}
Nếu tổng vật chất cộng chân không lớn hơn một, độ cong là dương. Nếu nhỏ hơn một, độ cong là âm.

Khi đó, logic quan sát là khá trực tiếp. Đo vật chất thông thường và vật chất tối. Đo sự giãn nở Hubble. Đo độ cong bằng hình học trên những thang rất lớn. Cũng lưu ý rằng vũ trụ đang gia tốc. Chỉ dấu của hạng chân không không phải là \(H\) tự nó mà là xu hướng của \(\dot a/a\) tiến tới một hằng số, điều đó có nghĩa là thừa số tỉ lệ tiến gần nhánh mũ,
\begin{equation}
a(t)\propto e^{Ht}.
\end{equation}
Các phép đo ấy kiểm tra chéo lẫn nhau và cho thấy vũ trụ gần như phẳng, trong sai khác chỉ vài phần trăm.

\begin{classroomqa}
\lecturetimestamp{00:10:50}
\audiencequestion{Vật chất có bao gồm vật chất tối không?}
\lecturerresponse{Có. Trong thảo luận hiện tại, hạng vật chất bao gồm cả vật chất tối. Cũng có thể thêm bức xạ, nhưng hiện nay phần đó nhỏ đến mức có thể bỏ qua.}
\end{classroomqa}

\begin{classroomqa}
\lecturetimestamp{00:11:11}
\audiencequestion{Có thước đo nào của năng lượng tối ngoài hằng số Hubble không?}
\lecturerresponse{Điều cốt yếu không phải là giá trị của \(H\) tự nó, mà là cách sự giãn nở tiến dần tới một định luật mũ. Cùng với phép đo vật chất và phép đo hình học của độ phẳng, điều đó cho phép suy ra hạng chân không.}
\end{classroomqa}

\section{Từ độ phẳng đến cấu trúc}

Vì sao vũ trụ lại phẳng đến thế? Lời giải thích được đề xuất là lạm phát. Phẳng có nghĩa là rất lớn, rất trơn nhẵn, và rất đồng nhất trên những thang cỡ hàng tỷ, đúng hơn là hàng chục tỷ, năm ánh sáng. Giãn nở mũ có thể kéo giãn một vùng cho đến khi nó trông như vậy.

Nhưng một điều tốt cũng có thể bị đẩy quá mức. Một vũ trụ hoàn toàn trơn nhẵn sẽ vẫn hoàn toàn trơn nhẵn. Hấp dẫn không tạo ra sự vón cục từ đối xứng chính xác. Nếu hiện nay có cấu trúc, thì trước đó nó phải bắt đầu từ một bất đồng nhất nhỏ. Lần ngược vũ trụ hiện nay về quá khứ, độ bất đồng nhất ở thời điểm trở nên trong suốt ra vào cỡ một phần trong \(10^5\):
\begin{equation}
\frac{\delta\rho}{\rho}\sim 10^{-5}.
\end{equation}
Vũ trụ đã cực kỳ trơn nhẵn, nhưng không hoàn toàn như thế.

Từ đó trở đi, hấp dẫn làm điều mà các lực thông thường trong một căn phòng đầy không khí không làm. Một vùng dư mật độ nhỏ kéo vật chất về phía nó. Vùng lân cận bị rút bớt vật chất, nên vùng thiếu mật độ càng thiếu hơn và vùng dư mật độ càng dư hơn. Hấp dẫn khuếch đại sự tương phản thay vì xóa nó đi.

Câu hỏi tự nhiên là vì sao điều này đã không xảy ra sớm hơn. Câu trả lời là áp suất. Trước đó, áp suất bức xạ lớn. Dùng
\begin{equation}
p=w\rho,
\end{equation}
các trường hợp quen thuộc là
\begin{equation}
w_{\rm rad}=\frac13,\qquad
w_m=0,\qquad
w_{\rm vac}=-1.
\end{equation}
Áp suất dương có tác dụng làm đồng nhất hóa. Nó đẩy các nhiễu loạn ra ngoài và làm trơn môi trường. Vật chất, khi được xem trên những thang lớn là không áp suất, không cung cấp sự chống đỡ như vậy. Khi có áp suất, hấp dẫn bị kìm lại. Khi áp suất trở nên không đáng kể, hấp dẫn thắng thế.

Vật chất tối giữ vai trò then chốt trong trình tự các sự kiện. Những khối sụp đổ đầu tiên là các quầng vật chất tối. Vật chất baryon thông thường sau đó rơi vào các quầng ấy và, vì baryon có thể tiêu tán, tập trung gần tâm hơn. Vật chất tối không tiêu tán theo cùng cách đó, nên nó vẫn phân bố trải khắp quầng. Sự bắt đầu của quá trình co đó sớm hơn thời điểm trở nên trong suốt một chút, với chênh lệch thời gian xấp xỉ một hệ số mười, đó là lý do phông vi sóng không cho thấy trực tiếp sự khởi phát ấy.

Một nhận xét ngắn làm rõ lịch sử hóa học về sau. Hầu như mọi nguyên tố nặng hơn heli đều là sản phẩm của tiến hóa sao chứ không phải của vũ trụ sơ khai. Vũ trụ sơ khai về cơ bản chỉ tạo ra hydro, heli, và một lượng vết nhỏ của vài hạt nhân nhẹ.

Cái được quan sát trên bầu trời không phải chính mật độ khối lượng mà là tính dị hướng nhiệt độ trong phông vi sóng. Các thăng giáng mật độ góp phần vào đó. Các thăng giáng vận tốc cũng góp phần. Các vùng dư mật độ có xu hướng trông lạnh hơn đôi chút vì photon mất thêm năng lượng khi leo ra khỏi các giếng thế hấp dẫn sâu hơn, trong khi các dịch Doppler do chuyển động của chất khí còn chồng thêm lên đó. Phân tích số đầy đủ thì phức tạp. Mệnh đề vững chắc thì đơn giản hơn: vào thời điểm vũ trụ trở nên trong suốt, mật độ năng lượng biến thiên cỡ một phần trong \(10^5\).

\begin{classroomqa}
\lecturetimestamp{00:20:29}
\audiencequestion{Vì sao hấp dẫn đã không làm các chỗ vón cục lớn lên sớm hơn?}
\lecturerresponse{Bởi vì trước đó vũ trụ bị bức xạ chi phối, và bức xạ mang đủ áp suất dương để làm đồng nhất môi trường. Chỉ khi bức xạ trở nên không còn quan trọng so với vật chất thì hấp dẫn mới bắt đầu chi phối sự tăng trưởng của các chỗ vón cục.}
\end{classroomqa}

\begin{classroomqa}
\lecturetimestamp{00:24:35}
\audiencequestion{Sự sụp đổ chẳng phải cần một cơ chế tiêu tán sao?}
\lecturerresponse{Vật chất baryon thông thường quả thật cần tiêu tán nếu muốn lắng về phía tâm. Nhưng vật chất tối có thể bắt đầu bằng cách tạo thành các quầng mà không cần kiểu tiêu tán do va chạm đó. Còn có một dạng suy giảm mang tính vũ trụ học gắn với chính sự giãn nở, và điều đó sẽ trở nên quan trọng về sau.}
\end{classroomqa}

\begin{classroomqa}
\lecturetimestamp{00:30:34}
\audiencequestion{Phải chăng khi đó cũng có các thăng giáng vận tốc, để về sau các thiên hà có thể thu được mômen động lượng?}
\lecturerresponse{Có. Mẫu hình sơ khai bao gồm các thăng giáng mật độ và các thăng giáng vận tốc. Cái được thấy trên bầu trời là một mẫu nhiệt độ được tạo ra bởi cả các biến thiên mật độ lẫn chuyển động của chất khí.}
\end{classroomqa}

\section{Inflaton và cao nguyên}

Cơ chế bây giờ được đưa vào là một trường vô hướng \(\phi\). Mật độ năng lượng của nó có ba phần:
\begin{equation}
\rho_\phi=\frac12\dot\phi^2+\frac12(\nabla\phi)^2+V(\phi).
\end{equation}
Biến thiên theo thời gian góp phần. Biến thiên theo không gian góp phần. Một trường vô hướng cũng có thể mang một thế năng thật sự \(V(\phi)\). Phần đó không có gì lạ; trường Higgs là ví dụ quen thuộc. Điều quan trọng ở đây là dạng của thế năng.

Thế năng phải khá phẳng trên một khoảng nào đó, với một độ dốc đi xuống nhẹ, rồi tại một giá trị tới hạn đáy đột ngột sụt xuống. Nó không rơi về đúng bằng 0. Nó rơi xuống một giá trị rất nhỏ. Theo đơn vị Planck, năng lượng chân không hiện nay có bậc độ lớn
\begin{equation}
10^{-123}.
\end{equation}
Ở phía trên cao nguyên, nó lớn hơn giá trị này rất nhiều, nhưng vẫn thấp hơn thang Planck rất xa.

\begin{figure}[t]
\centering
\includegraphics[width=0.78\textwidth]{lecture_08_figure_03.png}
\caption{Thế năng inflaton định tính với một cao nguyên rộng, một đoạn sụt dốc, và một cực tiểu thấp.}
\lectureframe{00:43:23}
\end{figure}

Đồ thị này không mô tả một vị trí trong không gian. Nó mô tả một giá trị của trường. Nói rằng vũ trụ bắt đầu ở cao trên cao nguyên chỉ có nghĩa là giá trị của trường, trên một vùng rất lớn, nằm ở đó. Ban đầu \(\phi\) có biến thiên đôi chút từ nơi này sang nơi khác cũng không quan trọng, vì lạm phát kéo giãn những biến thiên ấy lên rất lớn. Sau đủ nhiều giãn nở, bất kỳ vùng cục bộ nào cũng trông gần như đồng nhất. \editorialnote{Nhãn các trục không đọc được trong khung hình bảng, nhưng phần thảo luận bằng lời cho thấy rõ rằng trục đứng là thế năng \(V(\phi)\) còn trục ngang là giá trị trường \(\phi\).}

Chừng nào trường còn nằm trên phần rộng và dốc chậm của thế năng, thế năng còn hành xử như năng lượng chân không. Nếu \(\phi\) bị đóng băng hoàn toàn, nó sẽ hành xử đúng hệt như năng lượng chân không. Bỏ qua các hệ số số học đã lược đi, hệ thức được nêu trong bài giảng đơn giản là
\begin{equation}
H\propto \sqrt{V(\phi)},
\end{equation}
và điều đó cho định luật gần mũ
\begin{equation}
a(t)\propto e^{Ht}.
\end{equation}
Hầu như mọi thứ hiện nay còn có thể quan sát được đều xảy ra sau khi trường vượt qua mép đó. Chỗ sụt là nơi năng lượng lớn kiểu chân không trên cao nguyên được chuyển thành bức xạ, vật chất thông thường, vật chất tối, và phần năng lượng chân không còn dư rất nhỏ được thấy ngày nay.

\begin{classroomqa}
\lecturetimestamp{00:44:07}
\audiencequestion{Trường \(\phi\) là trường trên cái gì, và các phương trình Einstein có còn đang được dùng không?}
\lecturerresponse{Vâng. \(\phi\) là một trường trên không-thời gian, và vẫn dùng cùng phương trình Friedmann. Trong chế độ này, mật độ năng lượng ở vế phải chủ yếu do \(V(\phi)\) cung cấp, miễn là trường biến thiên đủ chậm.}
\end{classroomqa}

\begin{classroomqa}
\lecturetimestamp{00:47:04}
\audiencequestion{Làm thế nào năng lượng inflaton dạng chân không có thể biến thành vật chất và bức xạ?}
\lecturerresponse{Khi trường lăn xuống đáy và dao động, những dao động ấy chính là các lượng tử của trường. Nếu các lượng tử đó không bền, chúng có thể phân rã thành photon, electron, positron và các hạt quen thuộc khác.}
\end{classroomqa}

\section{Lăn chậm và ma sát Hubble}

Một khi trường vượt qua mép, vì sao nó lại không đơn giản cứ dao động mãi mãi? Vì sao nó lại không lắc qua lắc lại và có lẽ thậm chí còn leo ngược trở lên? Câu trả lời được đưa ra ở đây là một cơ chế tắt dần mang tính vũ trụ học. Nó không phải là ma sát vi mô thực sự, nhưng trong phương trình chuyển động nó hành xử như vậy.

Bỏ qua các gradient không gian và coi trường là đồng nhất. Khi đó hạng động năng và hạng thế năng trông giống như một bậc tự do duy nhất với động lực học \(T-V\), nhưng vì chúng là các mật độ năng lượng nên phải được nhân với thể tích vật lý \(a^3\). Từ đó ta đi tới một Lagrangian hiệu dụng đồng nhất có dạng
\begin{equation}
L_{\mathrm{hom}}=a^3(t)\left[\frac12\dot\phi^2-V(\phi)\right].
\end{equation}
\editorialnote{Dòng trên cùng trên bảng bị cắt mất một phần. Lagrangian đồng nhất này chỉ được phục dựng đến mức cần thiết để tái tạo phương trình chuyển động chắc chắn được viết bên dưới nó.}

Phương trình còn đọc chắc chắn được trên bảng là
\begin{equation}
\frac{d}{dt}\left(a^3\dot\phi\right)=-a^3\frac{\partial V}{\partial\phi}.
\end{equation}
Khai triển đạo hàm theo thời gian cho
\begin{equation}
a^3\ddot\phi+3a^2\dot a\,\dot\phi=-a^3\frac{\partial V}{\partial\phi},
\end{equation}
và chia cho \(a^3\) ta được
\begin{equation}
\ddot\phi+3H\dot\phi=-\frac{\partial V}{\partial\phi}.
\end{equation}
Dấu trừ được sửa một cách tường minh trong phần suy diễn bằng lời: lực là âm đạo hàm của thế năng.

\begin{figure}[t]
\centering
\includegraphics[width=0.78\textwidth]{lecture_08_figure_05.png}
\caption{Phương trình chuyển động đồng nhất của inflaton, cho thấy hạng ma sát Hubble \(3H\dot\phi\).}
\lectureframe{00:57:10}
\end{figure}

Hạng \(3H\dot\phi\) là cái được gọi là ma sát Hubble. Nó hoàn toàn xuất phát từ việc lấy đạo hàm thừa số \(a^3\). Tuy nhiên, về mặt động lực học, nó trông giống hạng lực cản trong chuyển động của một hòn đá rơi qua nước. Có một hạng lực, một hạng gia tốc và một hạng tuyến tính theo vận tốc. Trường nhanh chóng đạt tới vận tốc giới hạn rồi sau đó di chuyển chậm dọc theo cao nguyên. Đó là chế độ lăn chậm. Chừng nào độ dốc còn thoải và trường còn chuyển động chậm, thế năng chỉ thay đổi chậm, nên sự giãn nở vẫn gần như theo hàm mũ.

Gần cực tiểu, trường dao động. Những dao động ấy là các lượng tử inflaton. Nếu các lượng tử đó không bền, chúng phân rã, và đó là cách năng lượng trên cao nguyên được chuyển thành thành phần vật chất nóng và bức xạ của vũ trụ về sau.

\begin{classroomqa}
\lecturetimestamp{00:50:48}
\audiencequestion{Vì sao trường không đơn giản dao động mãi mãi, và cái gì tạo ra sự tắt dần?}
\lecturerresponse{Bản thân sự giãn nở tạo ra một hạng giống ma sát. Thừa số \(a^3\) trong Lagrangian đồng nhất sinh ra hạng \(3H\dot\phi\), và hạng đó hành xử như lực cản nhớt.}
\end{classroomqa}

\section{Thăng giáng lượng tử và sự vượt ra khỏi chân trời}

Giờ đây có vẻ nảy sinh một vấn đề. Nếu lạm phát làm vũ trụ trở nên cực kỳ trơn nhẵn, tại sao nó không xóa đi chính những mầm mà các thiên hà cần đến? Câu trả lời là cơ học lượng tử. Những thăng giáng còn lại không phải là tàn dư của một quá trình san phẳng không hoàn hảo. Chúng được tạo ra mới.

Tạm gác sự giãn nở sang một bên. Ngay cả khi một trường gần như đồng nhất hoàn toàn, cơ học lượng tử cũng không cho phép cả \(\phi\) lẫn \(\dot\phi\) đều được xác định chính xác. Trạng thái cơ bản vẫn bị khuấy động. Đây cũng là lý do khiến một dao động tử điều hòa thông thường có năng lượng điểm không. Trong các hoàn cảnh hằng ngày, những thăng giáng ấy khó thấy vì chúng có bước sóng ngắn và dao động nhanh, nên các quan sát trung bình hóa chúng.

Sự giãn nở làm thay đổi câu chuyện. Trong tọa độ đồng chuyển động, mode của trường giữ gần như cùng một bước sóng tọa độ, trong khi bước sóng vật lý tăng theo thừa số tỉ lệ:
\begin{equation}
\lambda_{\rm phys}=a(t)\,\Delta x.
\end{equation}
Trong thời kỳ lạm phát, thang Hubble vật lý xấp xỉ cố định,
\begin{equation}
d_H\sim H^{-1},
\end{equation}
nên chân trời đồng chuyển động tương ứng co lại theo
\begin{equation}
d_{H,\mathrm{coord}}\sim \frac{1}{aH}.
\end{equation}

Phép ví với sợi dây làm cho việc đóng băng trở nên trực quan. Một sóng dao động vì mỗi đoạn dây bị các đoạn lân cận kéo. Nhưng nếu bước sóng trở nên lớn bằng chân trời, thì những đoạn lân cận có liên quan đã vượt ra ngoài liên hệ nhân quả. Lực kéo phục hồi biến mất. Một khi mode dài hơn chân trời, nó bị đóng băng.

Điều này xảy ra lặp đi lặp lại. Các bước sóng dài hơn đóng băng trước. Các bước sóng ngắn hơn đóng băng sau. Ngay cả một mode bắt đầu với bước sóng rất nhỏ cuối cùng cũng bị kéo giãn cho đến khi nó vượt qua chân trời. Vì thế, lạm phát không đơn thuần xóa cấu trúc. Nó liên tục tạo ra những biến thiên bị đóng băng của trường.

Quá trình này cũng không dừng lại chỉ vì lạm phát kéo dài lâu. Các mode đã đóng băng từ trước bị kéo giãn tới những thang cực lớn, nhưng các mode mới luôn tiếp tục đóng băng. Mẫu hình tích lũy trở nên dừng về mặt thống kê theo nghĩa thích hợp cho vũ trụ học. Đó là nguồn gốc của phổ gần như bất biến theo tỉ lệ: cùng một quá trình vượt qua chân trời lặp lại trên các thang logarit kế tiếp nhau và tạo ra biên độ gần như như nhau từ dải bước sóng này sang dải bước sóng kế tiếp.

Sau khi trường lăn khỏi cao nguyên, chân trời lại lớn lên. Những mode trước đó ở bên ngoài chân trời lạm phát tạm thời về sau quay trở lại bên trong. Theo cách nói thông thường, chúng rời khỏi chân trời trong thời kỳ lạm phát và đi vào lại sau đó. Đến lúc ấy chúng đã bị kéo giãn tới các thang vĩ mô, nên khi bắt đầu dao động trở lại, chúng dao động dưới dạng những cấu trúc rất lớn.

\begin{classroomqa}
\lecturetimestamp{01:12:02}
\audiencequestion{Trên đồ thị không-thời gian, chân trời trông như thế nào?}
\lecturerresponse{Khoảng cách vật lý đến chân trời giữ gần như cố định ở \(c/H\) trong giai đoạn giãn nở theo hàm mũ, nhưng trên một đồ thị trong tọa độ đồng chuyển động, khoảng cách vật lý cố định ấy trông ngày càng nhỏ đi vì thang tọa độ đang bị kéo giãn. Theo ngôn ngữ đồng chuyển động, chân trời co lại như \(1/(aH)\).}
\end{classroomqa}

\begin{classroomqa}
\lecturetimestamp{01:15:59}
\audiencequestion{Biên độ biến đổi thế nào khi các thăng giáng bị kéo giãn, và có một trạng thái dừng hay không?}
\lecturerresponse{Các mode riêng lẻ bị suy giảm khi bị kéo giãn, nhưng các mode mới vẫn tiếp tục đóng băng. Kết quả là một mẫu hình dừng về mặt thống kê, gần như bất biến theo tỉ lệ, chứ không phải sự biến mất hoàn toàn của các thăng giáng.}
\end{classroomqa}

\section{Từ các biến thiên bị đóng băng của trường đến các biến thiên mật độ}

Một khi trường đã có được những biến thiên nhỏ từ nơi này sang nơi khác, những biến thiên ấy được chuyển thành chênh lệch mật độ khi các vùng khác nhau đi qua mép vào những thời điểm hơi khác nhau. Một vùng nằm thấp hơn một chút dọc sườn dốc sẽ tới vách trước. Nếu năng lượng chân không của nó được chuyển thành bức xạ, thì bức xạ ấy bị pha loãng theo sự giãn nở. Một vùng lân cận vẫn còn trên cao nguyên thì vẫn do chân không chi phối và không bị pha loãng. Đến khi vùng thứ hai rơi xuống, vùng thứ nhất đã bị dịch đỏ mất một phần năng lượng.

Đó là lý do vì sao sự lăn chậm hơn tạo ra độ tương phản mật độ cuối cùng lớn hơn. Nếu trường chỉ chuyển động rất chậm dọc theo cao nguyên, thì một khác biệt nhỏ về giá trị của trường tương ứng với một độ trễ thời gian lớn hơn giữa lúc vùng này và vùng kế tiếp chạm tới vách. Vùng rơi sớm hơn có nhiều thời gian hơn để bị pha loãng trước khi vùng rơi muộn hơn chuyển năng lượng của nó. Độ dốc càng thoải, vận tốc giới hạn càng nhỏ, và biến thiên mật độ cuối cùng càng lớn. Một cao nguyên cao hơn cũng có xu hướng cho các thăng giáng lớn hơn. Bề rộng của cao nguyên chủ yếu quyết định lạm phát kéo dài bao lâu.

Chỉ một tổ hợp của độ cao và độ dốc mới bị dữ liệu về thăng giáng ràng buộc chặt chẽ. Nếu thang năng lượng đủ cao, cũng phải có một đóng góp của sóng hấp dẫn có thể quan sát được. Các cận số chỉ được nêu như những ước lượng rất thô. Nói đại thể, nếu mật độ năng lượng trên cao nguyên lớn hơn nhiều so với khoảng \(10^{-14}\) tính theo đơn vị Planck, thì bức xạ nền vi sóng hẳn đã cho thấy quá nhiều cấu trúc do sóng hấp dẫn; còn nếu nhỏ hơn nhiều so với khoảng \(10^{-16}\), thì sẽ hầu như không có khả năng thực tế nào để thấy được một tín hiệu như vậy. \editorialnote{Người nói nêu rõ rằng các con số này chỉ là gần đúng và ông không nhớ các giá trị chính xác. Chúng được giữ lại ở đây chỉ như những phát biểu về bậc độ lớn.}

Các đơn vị được giải thích theo cách thông thường của hệ đơn vị tự nhiên. Bước sóng Compton của một hạt có khối lượng \(m\) là
\begin{equation}
\lambda_C=\frac{\hbar}{mc}.
\end{equation}
Nếu \(c=\hbar=1\), thì khối lượng và nghịch đảo độ dài có cùng đơn vị. Mật độ năng lượng là năng lượng trên một đơn vị thể tích, nên đơn vị của nó là khối lượng mũ bốn:
\begin{equation}
[\rho]=(\mathrm{mass})^4.
\end{equation}

Nhìn tại một điểm đơn lẻ, sự tích lũy của các mode bị đóng băng trông như một bước đi ngẫu nhiên. Các điểm khác nhau trong không gian trải qua những bước đi ngẫu nhiên khác nhau, và chính điều đó biến nhiễu lượng tử thành biến thiên theo không gian.

\begin{classroomqa}
\lecturetimestamp{01:33:04}
\audiencequestion{Liệu \(\phi\) có thực hiện một bước đi ngẫu nhiên không, và điều đó có khiến điểm khởi đầu từ sớm trở nên không còn quan trọng không?}
\lecturerresponse{Tại bất kỳ điểm cố định nào, các mode bị đóng băng kế tiếp nhau cộng lại với pha ngẫu nhiên, nên \(\phi\) quả thật thực hiện một bước đi ngẫu nhiên. Chính thành công đó cũng là khó khăn: lạm phát liên tục thay thế thông tin sớm hơn bằng các thăng giáng muộn hơn, nên giai đoạn trước lạm phát trở nên khó khôi phục.}
\end{classroomqa}

\section{Lạm phát che giấu điều gì}

Lý thuyết này thành công theo một nghĩa rất cụ thể. Nó liên hệ các dao động trên một dải thang độ khổng lồ với một cơ chế nền tảng duy nhất. Nhưng chính cơ chế đó lại xóa nhòa những gì có trước. Các bất quy tắc trước đó bị kéo giãn đến mức biến mất, và các dao động lượng tử được tạo ra về sau mới là những gì còn giữ vai trò liên quan. Đó là lý do lạm phát vừa hấp dẫn vừa gây bực bội. Nó giải thích tính trơn đều và phổ dao động, nhưng cũng xóa đi phần lớn lịch sử trước đó.

Điều này dẫn đến một sự băn khoăn về điểm xuất phát. Mật độ năng lượng rất cao không phải là vấn đề. Vấn đề là năng lượng của vùng bằng phẳng dường như không tự nhiên ở cỡ Planck, mà cũng không tự nhiên là cực nhỏ. Nó cao theo các đơn vị thông thường, nhưng vẫn nhỏ theo đơn vị Planck. Nó không có vẻ là một điểm dừng chân đầu tiên hiển nhiên cho một lý thuyết cơ bản.

Trong bức tranh đơn giản này cũng không có mốc \(t=0\) theo nghĩa đen. Một khi lạm phát kéo dài quá lâu, gần như mọi thông tin về bất kỳ giai đoạn nào trước đó đều biến mất. Nếu số e-folding chỉ vừa nhỉnh hơn mức tối thiểu cần thiết, có lẽ một vài dấu vết của giai đoạn trước vẫn có thể còn sót lại, nhưng nhìn chung lạm phát xóa sạch quá khứ.

Sự thận trọng tương tự cũng xuất hiện khi xét không gian de Sitter. Quy luật mũ thuần túy là đúng cho nhánh lạm phát ở thời muộn, nhưng nghiệm de Sitter chính xác được mô tả đúng hơn bởi
\begin{equation}
a(t)\propto \cosh(Ht).
\end{equation}
Nghiệm đó co lại đến một kích thước tối thiểu rồi giãn nở. Việc có nên hiểu nhánh co ấy theo nghĩa đen trong vũ trụ học hay không được để ngỏ.

\begin{classroomqa}
\lecturetimestamp{01:37:31}
\audiencequestion{Liệu có thể không hề có điểm bắt đầu nào, với hành vi mũ kéo dài vô hạn về quá khứ không?}
\lecturerresponse{Không trong bức tranh đơn giản này. Không gian de Sitter chính xác không phải là một hành vi mũ giãn nở vĩnh viễn trong mọi thời gian; nó giống \(\cosh(Ht)\) hơn, với một nhánh co và một nhánh giãn nở. Việc có nên hiểu nhánh sớm hơn ấy theo nghĩa đen hay không vẫn là một câu hỏi bỏ ngỏ.}
\end{classroomqa}

\section{Tinh chỉnh, độ cong và các cận nhân học}

Ở gần đúng bậc nhất, các tham số quan trọng là độ cao của vùng bằng phẳng, độ dốc của nó và độ rộng \(\Delta\phi\) của nó. Độ dốc quan trọng vì nó xác định vận tốc giới hạn và do đó xác định độ lớn của các dao động. Bề rộng quan trọng vì nó góp phần quyết định lạm phát kéo dài bao lâu. Một định kiến thường gặp là độ biến thiên cơ bản của trường không nên lớn hơn khối lượng Planck quá nhiều, và mật độ năng lượng không nên vượt quá cỡ một theo đơn vị Planck. Dưới những cận thô như vậy, miền cho đủ lạm phát không phủ kín toàn bộ không gian tham số. Ước lượng được nêu ở đây là cần tinh chỉnh ở mức vài phần trăm, khoảng \(2\%\), để làm cho độ dốc đủ thoải cho khoảng \(60\) đến \(62\) e-fold.

\begin{classroomqa}
\lecturetimestamp{01:46:02}
\audiencequestion{Có thể hiểu lạm phát theo ngôn ngữ không gian tham số không, và nó cần tinh chỉnh đến mức nào?}
\lecturerresponse{Có. Ở gần đúng bậc nhất, người ta nghĩ theo độ cao, độ dốc và bề rộng của vùng bằng phẳng. Bài giảng đưa ra ước lượng thô rằng để có đủ số e-fold thì cần tinh chỉnh ở mức vài phần trăm, khoảng \(2\%\).}
\end{classroomqa}

Mức tinh chỉnh tương đối nhẹ đó sau đó được đem so sánh với câu đố sắc nét hơn nhiều của hằng số vũ trụ học hiện nay, cỡ \(10^{-123}\). Phần thảo luận chuyển sang các câu hỏi nhân học vì cả hai vấn đề đều có thể được phát biểu lại bằng cách hỏi điều gì là cần thiết để cấu trúc hình thành.

Trong bức tranh Newton, một hằng số vũ trụ học dương tác dụng như một lực đẩy phổ quát tỉ lệ với khoảng cách. Nếu nó lớn hơn nhiều, để trở nên tai hại thì nó không cần phải phá tung các thiên hà ngày nay. Nó chỉ cần đủ lớn để ngăn các dao động nguyên thủy cỡ \(10^{-5}\) không bao giờ ngưng tụ thành các thiên hà. Ví dụ được nêu là một giá trị lớn hơn giá trị quan sát được khoảng hai bậc độ lớn đã đủ làm điều đó. Đây là điểm gắn với lập luận của Weinberg: hằng số vũ trụ học không cần chính xác bằng không; nó chỉ cần đủ nhỏ để cho phép cấu trúc hình thành.

Độ cong âm tác dụng theo cùng một chiều. Trong ngôn ngữ Newton của các bài giảng trước, \(k<0\) có nghĩa là sự lùi xa nhanh hơn vận tốc thoát. Lực đẩy hướng ra ngoài quá lớn cũng khiến các quá mật độ nguyên thủy nhỏ không bao giờ trở thành các thiên hà. Lạm phát pha loãng độ cong đó. Cận dưới được nêu vào khoảng \(59.5\) e-folding. Sau đó, một giá trị quanh \(62\) được lấy làm một lịch sử lạm phát tiêu biểu nằm an toàn phía trên ngưỡng.

Điều này làm thay đổi câu hỏi về xác suất. Nếu ta hỏi về thể tích thuần túy của không gian tham số cho đủ lạm phát, câu trả lời có thể chỉ là vài phần trăm. Nếu điều kiện hóa trên việc ở trong một vũ trụ nơi các thiên hà hình thành và có các quan sát viên tồn tại để nêu câu hỏi ấy, thì xác suất có điều kiện có thể trở nên lớn hơn nhiều. Phần thảo luận dừng lại ở đó: không phải với một lời giải cuối cùng, mà với những câu đố về tinh chỉnh, sự diễn giải lại chúng theo nhân học, và sự thừa nhận rằng đây là nơi các cuộc tranh luận lớn vẫn còn tồn tại.